%% 05-how-to-use.tex — Audience guide

\section{How to Use This Portfolio}

Different readers will want different papers.
This section points each audience to the right starting point.

\paragraph{If you are a judge or clerk.}
Start with this document and the Policy Brief (A.5).
The Policy Brief is two to four pages and gives you the legal and
empirical case in plain language, with citations to \emph{Rucho},
the Voting Rights Act, and the relevant papers.
Then read Paper D.1 (the 42\% VRA threshold) if the case involves
majority-minority districts, or Paper C.5 (the efficiency gap analysis)
if the case involves partisan fairness.
The quickstart guide at \texttt{docs/quickstart/quickstart-callais-expert.md}
covers Gingles bloc-voting methodology for expert witnesses.

\paragraph{If you are a state legislator or redistricting commissioner.}
Start with Paper B.1 (the algorithm overview) and the replication
package (A.4), which explains how to run the algorithm on any state
in under thirty minutes.
Paper D.4 addresses legal implementation pathways, including model
statutory language.
Paper F.1 covers state legislative districts if your chamber is also
being redistricted.
The quickstart guide at \texttt{docs/quickstart/quickstart-state-staff.md}
walks through the \texttt{redist state} command for your specific state.

\paragraph{If you are a researcher.}
The algorithm is documented in Track B (eighteen papers).
Reproducibility materials are in A.4; the full software stack,
SHA-256 audit chain, and Vermont walkthrough fixture are all included.
The ensemble diagnostics in Track G (especially G.4) provide
the statistical framework for comparing this algorithm against
Markov chain ensemble methods.
The quickstart guide at \texttt{docs/quickstart/quickstart-researcher.md}
covers the research API, Jupyter notebooks, and the
\texttt{redist analyze} command.

\paragraph{If you are a journalist.}
This document and the Policy Brief (A.5) are written for you.
The five headline numbers in Section~\ref{sec:findings} each carry
a paper citation for fact-checking.
The dashboard at \texttt{web/dashboard.html} lets you view maps
for any state and year in a browser without installing software.
For live contact, email is at the end of the Policy Brief.
